% \documentclass[draftcls, journal, onecolumn]{IEEEtran}
%\documentclass[journal,draftcls,onecolumn,12pt,twoside]{IEEEtranTCOM}
% \documentclass[journal,draftcls,onecolumn,12pt,twoside]{IEEEtranTCOM}
\documentclass[journal]{IEEEtran}
\usepackage{acronym}
%%%
\usepackage{bm}
\usepackage{booktabs}
% \usepackage{multirow}
% \usepackage{fancyhdr}
% \usepackage[cmex10]{amsmath}
% \usepackage{algorithm}
% \usepackage{algpseudocode}
% \usepackage{algorithmicx}
% \usepackage{filecontents}
% \usepackage[dvipsnames]{xcolor}
%\usepackage{mwe,tikz}
% \usepackage[percent]{overpic}
% \pagestyle{empty}
% \usepackage{todonotes}
\usepackage{tabularx}
% \usepackage{arydshln}
% \usepackage{verbatim}
%\usepackage{ulem} % sout
% \usepackage{gensymb}
\usepackage{stfloats}% <-- added
% \usepackage{mathrsfs}
%%%
\usepackage{siunitx}
\DeclareSIUnit{\atm}{atm}
\DeclareSIUnit{\bar}{bar}
\normalsize
\usepackage{amsmath,amssymb,amsfonts}
\usepackage{amsthm}
\newtheorem{lemma}{Lemma}
% \usepackage{algorithmic}
% \usepackage{graphicx, nicefrac}
% \usepackage{textcomp}
% \usepackage{xcolor}
% \usepackage{subcaption}
% \usepackage[ruled,vlined]{algorithm2e}
% \usepackage{ulem}
\usepackage{setspace} 
\usepackage{mathtools}
%vv% -- Packages from Infocom
%\usepackage[cmex10]{amsmath}
%\usepackage{amssymb}
\usepackage{color}
\usepackage{array}
\usepackage[version=4]{mhchem}
%\usepackage{graphicx}
%\usepackage{subfig}
\usepackage{float}
\usepackage{stmaryrd}
%\usepackage{enumitem}
\usepackage{soul}
\usepackage{stackrel}
\usepackage{epstopdf}
\usepackage{url}
% 
%^^%
\usepackage{cite}
%\usepackage[final]{graphicx}
% \usepackage{subfig}
% \usepackage{enumitem}
\usepackage{tikz}
%\usetikzlibrary{shapes,arrows,positioning}
\usetikzlibrary{automata,arrows,positioning,calc}
\usepackage{bigints}
% \usepackage{lineno} % To obtain line numbers
\usepackage{bbm}
 \usepackage{dsfont}
 \newcolumntype{L}[1]{>{\raggedright\arraybackslash}m{#1}} % Left-aligned with vertical centering
\newcolumntype{C}[1]{>{\centering\arraybackslash}m{#1}}  
% *** GRAPHICS RELATED PACKAGES ***
%
\ifCLASSINFOpdf
  % \usepackage[pdftex]{graphicx}
  % declare the path(s) where your graphic files are
  % \graphicspath{{../pdf/}{../jpeg/}}
  % and their extensions so you won't have to specify these with
  % every instance of \includegraphics
  % \DeclareGraphicsExtensions{.pdf,.jpeg,.png}
\else
  % or other class option (dvipsone, dvipdf, if not using dvips). graphicx
  % will default to the driver specified in the system graphics.cfg if no
  % driver is specified.
  % \usepackage[dvips]{graphicx}
  % declare the path(s) where your graphic files are
  % \graphicspath{{../eps/}}
  % and their extensions so you won't have to specify these with
  % every instance of \includegraphics
  % \DeclareGraphicsExtensions{.eps}
\fi


%\hyphenation{op-tical net-works semi-conduc-tor}

\acrodef{1D}{one-dimensional}
\acrodef{2D}{two-dimensional}
\acrodef{3D}{three-dimensional}
\acrodef{MC}{molecular communication}
\acrodef{TX}{transmitter}
\acrodef{RX}{receiver}
\acrodef{3D}{three dimensional}
\acrodef{CIR}{channel impulse response}
\acrodef{PDE}{partial differential equation}
\acrodef{ODE}{ordinary differential equation}

\begin{document}

\title{Title}
%
\author{
    Johannes Konrad,
    Maximilian Schäfer, and
    Fardad Vakilipoor
    \\
    % \IEEEauthorrefmark{1}Friedrich-Alexander-Universität Erlangen-Nürnberg, Erlangen, Germany\\
    % \IEEEauthorrefmark{2}Hochschule Geisenheim University, Geisenheim, Germany
\thanks{This paper was presented in part at the 13th ACM International Conference on Nanoscale Computing and Communication(cite the confi here) [DOI: add doi later.}
\thanks{Johannes Konrad,
    Maximilian Sch\"afer, and
    Fardad Vakilipoor are with the Institute for Digital Communications, Friedrich-Alexander-Universität Erlangen-Nürnberg, Germany (e-mail: johannes.konrad@fau.de).}}

% and the end of the author line. i.e., if you had this:
% 
% \author{....lastname \thanks{...} \thanks{...} }
%                     ^------------^------------^----Do not want these spaces!
%
% a space would be appended to the last name and could cause every name on that
% line to be shifted left slightly. This is one of those "LaTeX things". For
% instance, "\textbf{A} \textbf{B}" will typeset as "A B" not "AB". To get
% "AB" then you have to do: "\textbf{A}\textbf{B}"
% \thanks is no different in this regard, so shield the last } of each \thanks
% that ends a line with a % and do not let a space in before the next \thanks.
% Spaces after \IEEEmembership other than the last one are OK (and needed) as
% you are supposed to have spaces between the names. For what it is worth,
% this is a minor point as most people would not even notice if the said evil
% space somehow managed to creep in.


% The paper headers
\markboth{Konrad \MakeLowercase{et al.}: title (somehow abbreviated terms, or shorter version}%
{Konrad \MakeLowercase{\emph{et al.}}}

% make the title area
\maketitle

% As a general rule, do not put math, special symbols or citations
% in the abstract or keywords.
\begin{abstract}
%
Abstract...
\end{abstract}
%%
%%
\begin{IEEEkeywords}
molecular communication, agriculture, pest control.
\end{IEEEkeywords}
%%
%%
\acresetall
\IEEEpeerreviewmaketitle
%%
\section{Introduction}
Introduce \\
Motivate \\
Literature Review\\
List Contributions (bullet points)\\
Outline of the paper\\
%
% The remainder of the paper is organized as follows. Section~\ref{sec:agriMC} provides background on agricultural \ac{MC}; Section~\ref{sec:hipvMC} introduces the \ac{HIPV}-based \ac{MC} system and \ac{TX} design; Section~\ref{sec:system} derives the Reynolds-averaged advection-diffusion channel model; Section~\ref{sec:Analytical} presents its closed-form analytical solution under both ground boundary conditions; Section~\ref{sec:CEI} introduces the \ac{CEI} coverage metric; Section~\ref{sec:simulation} presents simulation results across \ac{TX} arrangements, release times, dates, and locations; and Section~\ref{sec:conclusion} concludes the paper.
%%
%%
\section{Section 1}

\section{Section 2}

\section{Section 3}
%%
%%
\section{Conclusion}\label{sec:conclusion}

%
% \appendices

\ifCLASSOPTIONcaptionsoff
  \newpage
\fi
\bibliographystyle{IEEEtran}
\bibliography{Bibliography}


\end{document}